O Memória de Curto e Longo Prazo (Long Short-Term Memory — LSTM) é uma variante das Redes Neurais Recorrentes (Recurrent
Neural Networks — RNN) projetada para resolver o problema do gradiente desaparecendo (vanishing gradient problem) que
afeta as RNNs quando precisam aprender dependências de longo prazo em sequências extensas. Enquanto uma RNN padrão
transmite apenas seu estado oculto (hidden state) de um passo ao outro, sofrendo perda de informações importantes ao
longo do tempo, o LSTM introduz um estado interno ou célula de memória que pode reter informações relevantes por vários
\textit{timesteps}, de acordo com sua importância.

Cada célula de LSTM é composta por quatro camadas neurais distintas: três portões — portão de esquecimento (forget
gate), portão de entrada (input gate) e portão de saída (output gate) — e uma camada responsável por gerenciar o próprio
estado interno. Esses portões combinam funções sigmoide e multiplicações ponto a ponto para controlar de forma
diferenciada o que deve ser descartado, adicionado ou propagado ao passo seguinte. A figura a seguir compara o fluxo de
informação em uma RNN simples e em uma célula de LSTM.  
% FIGURE: insira “LSTM3-SimpleRNN.png” e “LSTM3-chain.png” para comparação RNN × LSTM

O portão de esquecimento (forget gate) decide quais componentes de \(C_{t-1}\), o estado interno anterior, serão
removidos:
\[
f_t = \sigma\bigl(W_f[h_{t-1}, x_t] + b_f\bigr),
\]
onde \(\sigma\) é a sigmoide, \(h_{t-1}\) o estado oculto anterior, \(x_t\) a entrada atual, e \(W_f, b_f\) pesos e
vieses. Valores de \(f_t\) próximos a 0 indicam descarte, e próximos a 1 indicam retenção.  
% FIGURE: insira “LSTM3-focus-f.png” ilustrando o portão de esquecimento

O portão de entrada (input gate) regula a incorporação de novas informações através de duas etapas:  
1. Um cálculo sigmoide determina as posições a serem atualizadas:
   \[
   i_t = \sigma\bigl(W_i[h_{t-1}, x_t] + b_i\bigr).
   \]
2. Uma camada \(\tanh\) gera o vetor de candidatos:
   \[
   \widetilde{C}_t = \tanh\bigl(W_c[h_{t-1}, x_t] + b_c\bigr),
   \]
   com valores em \([-1,1]\).  
% FIGURE: insira “LSTM3-focus-i.png” ilustrando o portão de entrada

Em seguida, o estado interno é atualizado por
\[
C_t = f_t \,\odot\, C_{t-1} \;+\; i_t \,\odot\, \widetilde{C}_t,
\]
onde \(\odot\) denota multiplicação ponto a ponto.  
% FIGURE: insira “LSTM3-focus-C.png” mostrando a atualização de \(C_t\)

Finalmente, o portão de saída (output gate) define a porção de \(C_t\) que se tornará o próximo estado oculto:
\[
o_t = \sigma\bigl(W_o[h_{t-1}, x_t] + b_o\bigr),
\qquad
h_t = o_t \,\odot\, \tanh(C_t).
\]
% FIGURE: insira “LSTM3-focus-o.png” ilustrando o portão de saída

As equações essenciais de um LSTM resumem-se a:
\[
\begin{aligned}
f_t &= \sigma(W_f[h_{t-1}, x_t] + b_f),\\
i_t &= \sigma(W_i[h_{t-1}, x_t] + b_i),\\
\widetilde{C}_t &= \tanh(W_c[h_{t-1}, x_t] + b_c),\\
C_t &= f_t \,\odot\, C_{t-1} + i_t \,\odot\, \widetilde{C}_t,\\
o_t &= \sigma(W_o[h_{t-1}, x_t] + b_o),\\
h_t &= o_t \,\odot\, \tanh(C_t).
\end{aligned}
\]

A função de ativação \(\tanh\) pode ser substituída por outras, como ReLU, conforme a aplicação. Graças à combinação de
portões e do estado interno, o LSTM retém informações críticas por longos intervalos e descarta dados irrelevantes,
superando com eficiência o problema do gradiente desaparecendo e viabilizando o aprendizado de dependências de longo
prazo em tarefas de séries temporais e processamento de linguagem natural.



% % \section*{Referências} \begin{itemize} \item \url{https://colah.github.io/posts/2015-08-Understanding-LSTMs/} \item
% % \url{https://d2l.ai/chapter_recurrent-modern/lstm.html} \item %
% \url{https://www.tensorflow.org/guide/keras/sequential_model} \end{itemize}
